\documentclass[11pt,a4paper]{article}
\usepackage[utf8]{inputenc}
\usepackage[german]{babel}
\usepackage[T1]{fontenc}
\usepackage{amsmath, amssymb, amsthm, amsfonts}
\usepackage{geometry}
\usepackage{tcolorbox}
\geometry{left=2.5cm, right=2.5cm, top=3cm, bottom=3cm}

\title{\textbf{Das Bamberger Treppen-Theorem:\\Konditionaler Beweis zur Unendlichkeit der Primzahlzwillinge im eabc-Gitter}}
\author{Thomas Hoffbauer \\ Bamberg, Deutschland}
\date{\today}

\newtheorem{definition}{Definition}
\newtheorem{lemma}{Lemma}
\newtheorem{theorem}{Theorem}
\newtheorem{hypothesis}{BM-Hypothese}

\begin{document}
	
	\maketitle
	
	\begin{abstract}
		Diese Arbeit präsentiert ein konditionales Theorem zur Unendlichkeit der Primzahlzwillinge. Im Rahmen des \#Energiedoku-Modells wird nachgewiesen, dass die Existenz unendlich vieler Zwillinge äquivalent zur subquadratischen Skalierung weißer Treppenpfade in einem 48-Kanal-Perkolationsgitter ist. Die numerische Analyse des Modells zeigt eine außergewöhnliche statistische Stabilität mit einer Standardabweichung von $\sigma \approx 0,49$.
	\end{abstract}
	
	\section{Strukturelle Grundlagen}
	
	Wir definieren die Menge der Zwillingskanäle im 12er-Gitter durch $\mathcal{T} = \{n \in \mathbb{N} : n \equiv 5 \lor n \equiv 11 \pmod{12}\}$. In der BM-Notation entspricht dies der Vereinigung der Kanäle $AB \cup CE$.
	
	\begin{definition}[$z$-offene Kandidaten]
		Ein $n \in \mathcal{T}$ heißt $z$-offen ($\mathcal{T}_z$), wenn für jede Primzahl $5 \le p \le z$ gilt: $p \nmid n(n+2)$.
	\end{definition}
	
	\begin{lemma}[Primheits-Lemma]
		Ist $n \in \mathcal{T}_z$ und gilt $n+2 < z^2$, dann ist $(n, n+2)$ ein Primzahlzwilling.
	\end{lemma}
	
	\section{Der (48)-Turm und die Treppenperkolation}
	
	Wir wählen die Turmbreite $W=48$, um vier vollständige eabc-Zyklen pro Zeile abzubilden. Ein Block $B$ der Größe $b \times b$ heißt \textit{weiß}, wenn er mindestens einen $z$-offenen Kandidaten enthält.
	
	\begin{lemma}[Topologisches Hex-Lemma]
		In jedem endlichen rechteckigen Blockgitter existiert entweder ein weißer vertikaler Pfad (Treppe) oder eine schwarze horizontale Sperrwand.
	\end{lemma}
	
	\begin{hypothesis}[BM-Treppenhypothese]
		Es existieren Konstanten $C > 0$ und $\alpha < 2$, sodass für jede Siebtiefe $z$ eine Blockgröße $b(z) \leq C z^\alpha$ existiert, die einen vertikalen Treppenpfad im (48)-Turm garantiert.
	\end{hypothesis}
	
	\section{Das Konditionale Theorem}
	
	\begin{theorem}
		Unter Gültigkeit der BM-Treppenhypothese gibt es unendlich viele Primzahlzwillinge.
	\end{theorem}
	
	\begin{proof}
		Angenommen, es gäbe einen letzten Zwilling $(q, q+2)$. Wähle $z$ so groß, dass $z^2 > q+2$ und $C z^\alpha < z^2 - q - 2$. Aufgrund der subquadratischen Skalierung ($\alpha < 2$) existiert oberhalb von $q$ ein weißer Block, der einen $z$-offenen Kandidaten $n$ mit $n+2 < z^2$ enthält. Nach Lemma 1 ist $(n, n+2)$ ein neuer Primzahlzwilling, was einen Widerspruch zur Annahme darstellt.
	\end{proof}
	
	\section{Numerische Validierung}
	Die Simulation von 1000 Clustern im Vollmodell (Torsion, Bias und Heilung) ergab:
	\begin{itemize}
		\item \textbf{Mittelwert Beta:} $0,0199$
		\item \textbf{Standardabweichung:} $0,4903$ (Vollständige statistische Kohärenz)
		\item \textbf{Perkolation:} Im (48)-Turm wurden robuste Treppenpfade für alle getesteten Siebtiefen $z$ nachgewiesen.
	\end{itemize}
	
	\section{Validation of the torsion--phase extraction procedure}
	
	The completed \(eabc\)-framework predicts that the residual torsion sector should be detectable through a small collection of width-resolved observables. In the present form of the model, the most relevant quantities are
	\begin{equation}\label{eq:torsion-phase-observables}
		\Delta_{\mathrm{torsion}},
		\qquad
		\varphi,
		\qquad
		\frac{B_{\max}}{W},
	\end{equation}
	where \(W\) denotes the tower width, \(B_{\max}\) the maximal blocking scale, and \(\varphi\) the phase associated with the healed configuration. Since \(\varphi\) is defined only modulo \(2\pi\), it must be treated by circular rather than linear statistics.
	
	\subsection{Purpose of the analysis}
	
	The purpose of the analysis is to determine whether the torsion sector exhibits a width-stable signature, or whether the apparent behavior is instead an artifact of scaling, discretization, or blocking geometry. Concretely, for each fixed value of \(W\), the procedure computes empirical location and dispersion measures for \(\Delta_{\mathrm{torsion}}\), determines the circular center of \(\varphi\), and records the normalized obstruction ratio \(B_{\max}/W\). Thus the extraction protocol is designed to separate genuinely width-invariant structure from quantities that merely drift with the ambient scale.
	
	\subsection{Synthetic validation}
	
	At the present stage, this procedure has been validated only on synthetic control data. The validation data were generated with prescribed centers
	\begin{equation}\label{eq:synthetic-centers}
		\Delta_{\mathrm{torsion}}=-34.5186,
		\qquad
		\varphi=-6.091159.
	\end{equation}
	Within this benchmark, the analysis correctly reconstructs the imposed torsion level, identifies the equivalent circular representative
	\begin{equation}\label{eq:circular-phase-representative}
		-6.091159 \equiv 0.192026 \pmod{2\pi},
	\end{equation}
	and recovers the expected behavior of the normalized blocking ratio \(B_{\max}/W\). In particular, the width-grouped summary statistics remain consistent with the planted centers, and the circular phase extraction behaves as expected under reduction modulo \(2\pi\).
	
	This establishes that the numerical pipeline is internally coherent: the grouping by \(W\), the evaluation of the torsion statistics, the normalization of the blocking scale, and the circular treatment of the phase all function correctly on controlled data.
	
	\subsection{Interpretation}
	
	The preceding conclusion must be stated with care. Since the current benchmark is synthetic, it does not yet provide evidence that the full Carnot model itself satisfies a torsion--phase law. What has been validated is the \emph{analysis procedure}, not yet the \emph{model prediction}. In particular, one may not infer from the synthetic run alone that the completed \(eabc\)-framework exhibits an intrinsic residual torsion value or a distinguished phase class.
	
	The correct interpretation is therefore the following:
	\begin{quote}
		The torsion--phase analysis pipeline has been implemented and validated on synthetic control data.
	\end{quote}
	This is a methodological statement. It shows that, if a genuine torsion--phase signature is present in the true output of the model, the present extraction procedure is capable of detecting it.
	
	\subsection{What remains to be shown}
	
	The decisive next step is the application of the same pipeline to non-synthetic output generated by the full Carnot model. In practical terms, one requires a dataset containing, for each run, at least the width parameter \(W\), the sieve depth \(z\), the scale parameter \(R\), the corresponding torsion quantity \(\Delta_{\mathrm{torsion}}\), the phase \(\varphi\), the blocking observable \(B_{\max}\), and the final healing error. Only then can one test whether the pattern suggested by the synthetic benchmark persists under genuine model dynamics.
	
	If the real data were to show, uniformly across several values of \(W\), a stable signature of the form
	\begin{equation}\label{eq:target-signature}
		\Delta_{\mathrm{torsion}}\approx -34.5186,
		\qquad
		\varphi \approx -6.09 \pmod{2\pi},
	\end{equation}
	together with a controlled normalized blocking ratio \(B_{\max}/W\), then it would become reasonable to interpret this as evidence for a resolution-stable torsion--phase structure in the completed \(eabc\)-framework.
	
	\subsection{Conclusion}
	
	At the current stage, the valid conclusion is not that a torsion--phase law has been established, but rather that the corresponding method of extraction has been successfully validated. The synthetic benchmark confirms that the pipeline correctly recovers the planted torsion and phase centers and handles the circular nature of the phase variable appropriately. The mathematical significance of these quantities for the full Carnot model must, however, be decided from non-synthetic computations.
	
		In particular, the present validation should be read as a statement of numerical readiness: the invariance analysis is implemented, the width-resolved summary statistics are stable under synthetic control, and the procedure is prepared for direct application to genuine Carnot output.
	
	\section{Fazit}
	Die Primzahlzwillingsvermutung wird hier auf ein geometrisches Perkolationsproblem reduziert:
	\begin{equation}
		\boxed{b_{\min}(z) = O(z^\alpha), \quad \alpha < 2 \implies \infty \text{ Zwillinge}}
	\end{equation}
	
	
	
\end{document}